% Solutions/hints for ch02 exercises (assembled into Appendix C)

\begin{solution}{exr:ch02-time-expanded}
The time-expanded graph with horizon $T$ has $\abs{V}\,(T+1)$ vertices, one copy of every vertex per time layer. Between two consecutive layers there are $\abs{V}$ wait edges and, for an undirected base graph, $2\abs{E}$ move edges (each edge $\set{u,v}$ can be traversed in both directions), so there are $T\,(\abs{V}+2\abs{E})$ edges in total. Every edge goes from layer $t$ to layer $t+1$, so no edge ever returns to an earlier layer and the graph is acyclic. A time-indexed path $(v_0,\dots,v_T)$ visits one vertex per layer, and consecutive vertices are either equal (a wait edge) or adjacent in $G$ (a move edge); hence it is exactly a path from $(v_0,0)$ to $(v_T,T)$ in the time-expanded graph.
\end{solution}

\begin{solution}{exr:ch02-costs}
For the plan of \cref{ex:ch02-costs} the times are $T_1=3$ and $T_2=3$, so $\makespan=3$ and $\sumcost=6$. For the instance in which the two objectives disagree, take a corridor of cells $(0,1),\dots,(5,1)$ with agent~$A$ going from $(0,1)$ to $(5,1)$, agent~$B_1$ from $(1,2)$ to $(1,0)$ and agent~$B_2$ from $(2,3)$ through $(2,2)$ to $(2,0)$; every other cell is blocked. Without waiting, $B_1$ would enter cell $(1,1)$ at $t=1$ and $B_2$ would enter $(2,1)$ at $t=2$, exactly when $A$ is there. If $A$ never waits ($T_A=5$), both crossing agents must wait one step: $T_{B_1}=3$, $T_{B_2}=4$, giving $\makespan=5$ and $\sumcost=12$. If instead $A$ waits one step at its start, neither crossing agent is delayed: $T_A=6$, $T_{B_1}=2$, $T_{B_2}=3$, giving $\makespan=6$ and $\sumcost=11$. No plan does better on both counts. Every agent needs at least its unobstructed travel time ($5$, $2$ and $3$ steps), so $\sumcost\ge10$ with equality only if nobody waits, which is impossible because of the two conflicts. Hence $\sumcost\ge11$, and a plan with $\sumcost=11$ contains exactly one unit of delay; a delay of $B_1$ alone leaves the conflict of $A$ with $B_2$, a delay of $B_2$ alone leaves the conflict with $B_1$, so the delayed agent must be $A$ and the makespan is $6$. Conversely, makespan $5$ forces $A$ to move without waiting, so $B_1$ cannot enter $(1,1)$ before $t=2$ and $B_2$ cannot enter $(2,1)$ before $t=3$, which gives $\sumcost\ge5+3+4=12$. The two objectives therefore disagree on this instance.
\end{solution}

\begin{solution}{exr:ch02-closest-approach}
With $\pos=\pos_B-\pos_A$ and $\vel=\vel_B-\vel_A$ the squared distance is $D(t)^2=\norm{\pos+t\vel}^2=\pos\cdot\pos+2t\,\pos\cdot\vel+t^2\,\vel\cdot\vel$, a parabola in $t$. Differentiating the square gives $\tfrac{d}{dt}D(t)^2=2\,\pos\cdot\vel+2t\,\vel\cdot\vel$, which vanishes at $t^{*}$; solving for $t^{*}$ gives $t^*=-(\pos\cdot\vel)/(\vel\cdot\vel)$. Writing $\pos=\pos_\parallel+\pos_\perp$ with $\pos_\parallel$ along $\vel$, the term $\pos_\parallel+t^*\vel$ vanishes, so $d_{\min}=\norm{\pos_\perp}=\abs{\pos\times\vel}/\norm{\vel}$ because $\abs{\pos\times\vel}=\norm{\pos}\,\norm{\vel}\,\abs{\sin\phi}$ is the area of the parallelogram spanned by the two vectors. For $\pos=(4,2)$, $\vel=(-1,-1)$: $t^*=6/2=3$, $\pos\times\vel=-4+2=-2$, $d_{\min}=2/\sqrt{2}=\sqrt{2}$. Along the line of motion the disc of radius $R$ is entered at $t_c=t^*-\sqrt{R^2-d_{\min}^2}/\norm{\vel}$, which for $R=1.5$ gives $t_c=3-\sqrt{0.25}/\sqrt{2}\approx 2.646$; for $R=1<\sqrt{2}$ the square root is undefined and the discs never touch, in agreement with \code{time\_to\_collision} in \code{ch02\_toolbox.py}.
\end{solution}

\begin{solution}{exr:ch02-ellipse}
The characteristic polynomial of $\mat{\Sigma}=\begin{psmallmatrix}2&1.2\\1.2&1\end{psmallmatrix}$ is $\lambda^2-3\lambda+0.56=0$, so $\lambda_1=2.8$ and $\lambda_2=0.2$. The eigenvector of $\lambda_1$ solves $-0.8\,e_x+1.2\,e_y=0$, hence $\vect{e}_1\propto(3,2)$ and the major axis makes the angle $\arctan(2/3)\approx 33.7^\circ$ with the $x$-axis. The $1\sigma$ semi-axes are $\sqrt{2.8}\approx 1.673$ and $\sqrt{0.2}\approx 0.447$; the $2\sigma$ ellipse doubles both. In two dimensions the mass inside the $k\sigma$ ellipse is $1-e^{-k^2/2}$: $39.3\%$ for $k=1$ and $86.5\%$ for $k=2$ (not $68\%$ and $95\%$, which are the one-dimensional values); the $400$ seeded samples of \cref{fig:ch02-gaussian} give $39.8\%$ and $85.8\%$.
\end{solution}

\begin{solution}{exr:ch02-adjacency}
With the map \code{["....", ".\#..", "...."]} the blocked cell is $(1,1)$. Cell $(0,0)$ has the neighbours $(1,0)$ and $(0,1)$ at cost $1$; the diagonal to $(1,1)$ is blocked. Cell $(2,1)$ has $(1,1)$ blocked, so its neighbours are $(3,1)$, $(2,0)$, $(2,2)$ at cost $1$ and $(1,0)$, $(3,0)$, $(1,2)$, $(3,2)$ at cost $\sqrt2$ (the diagonals to $(1,0)$ and $(1,2)$ pass between $(1,1)$ and a free cell, and the corner-cutting rule of \cref{def:ch02-grid} forbids them, so only $(3,0)$ and $(3,2)$ remain). Cell $(3,0)$ has $(2,0)$ and $(3,1)$ at cost $1$ and $(2,1)$ at cost $\sqrt2$. For the counts: a 4-connected $W\times H$ grid has $W(H-1)$ vertical and $H(W-1)$ horizontal edges, and each of the $(W-1)(H-1)$ unit squares contributes two diagonals. For $W=4$, $H=3$ this gives $4\cdot2+3\cdot3=17$ and $17+2\cdot3\cdot2=29$.
\end{solution}

\begin{solution}{exr:ch02-inflation}
A cell centre at integer offset $(i,j)$ from the blocked cell is at distance $\max(\abs{i}-\tfrac12,0)$ and $\max(\abs{j}-\tfrac12,0)$ from the obstacle square in the two axes, so the distance to the square is $\sqrt{\max(\abs{i}-\frac12,0)^2+\max(\abs{j}-\frac12,0)^2}$. The thresholds are therefore $0.5$ (the four side neighbours), $\sqrt{0.5}\approx0.707$ (the four diagonal neighbours), $1.5$, $\sqrt{2.5}\approx1.581$ and $1.5\sqrt2\approx2.121$. Hence $r=0.6$ blocks $5$ cells, $r=1.0$ blocks $9$, and $r=1.6$ blocks $21$: the $3\times3$ block, the four cells at distance $1.5$ and the eight at distance $\sqrt{2.5}$. The blocked-cell count is a sampled area, and the exact Minkowski sum of the unit square with the disc of radius $r$ has area $1+4r+\pi r^2$; the ratio of the two tends to $1$ as $r$ grows, because the sampling error lives on the boundary and grows only like $r$.
\end{solution}

\begin{solution}{exr:ch02-stopping}
Braking from $v=2$~m/s with $a_{\max}=1$~m/s$^2$ takes $v/a_{\max}=2$~s, that is $20$ steps of $\dt=0.1$~s. The exact model \cref{eq:ch02-double-integrator} covers $v^2/(2a_{\max})=2.0$~m. Forward Euler updates the position with the \emph{old} velocity, so it adds $\dt\sum_{k=0}^{19}(2-0.1k)=0.1\,(2\cdot20-0.1\cdot190)=2.1$~m: it overestimates by $\tfrac12\dt^2 a_{\max}$ per step, $20\cdot0.005=0.1$~m in total. The admissibility rule is $v\le\sqrt{2a_{\max}d}$, equivalently $v^2/(2a_{\max})\le d$; \cref{ch:ch14} uses it to prune the dynamic window.
\end{solution}

\begin{solution}{exr:ch02-tangents}
Here $d=4$, $r=2$, so $\cos\alpha=r/d=\tfrac12$, $\alpha=60^\circ$ and $\sin\alpha=\sqrt3/2$. With $\vect{u}=(\vect{p}-\vect{c})/d=(-1,0)$ and $\vect{u}^{\perp}=(0,-1)$, \cref{eq:ch02-tangent-points} gives $\vect{t}_{\pm}=(4,0)+2\bigl(\tfrac12(-1,0)\pm\tfrac{\sqrt3}{2}(0,-1)\bigr)=(3,\mp\sqrt3)$. The tangent length is $\ell=\sqrt{d^2-r^2}=\sqrt{12}=2\sqrt3$ and the half-angle is $\theta=\arcsin(r/d)=30^\circ$. In general the right triangle $\vect{p}\vect{t}\vect{c}$ has the right angle at $\vect{t}$, hypotenuse $d$ and opposite leg $r$, so $\sin\theta=r/d$. As $d\to r^{+}$ the cone opens to a half-plane ($\theta\to90^\circ$) and as $d\to\infty$ it closes to a single direction. For a drone this says that an intruder that is already close forbids almost every velocity, which is why avoidance must start while $d$ is still large.
\end{solution}
